% Options for packages loaded elsewhere
\PassOptionsToPackage{unicode}{hyperref}
\PassOptionsToPackage{hyphens}{url}
%
\documentclass[
  10pt,
]{article}
\usepackage{amsmath,amssymb}
\usepackage{iftex}
\ifPDFTeX
  \usepackage[T1]{fontenc}
  \usepackage[utf8]{inputenc}
  \usepackage{textcomp} % provide euro and other symbols
\else % if luatex or xetex
  \usepackage{unicode-math} % this also loads fontspec
  \defaultfontfeatures{Scale=MatchLowercase}
  \defaultfontfeatures[\rmfamily]{Ligatures=TeX,Scale=1}
\fi
\usepackage{lmodern}
\ifPDFTeX\else
  % xetex/luatex font selection
\fi
% Use upquote if available, for straight quotes in verbatim environments
\IfFileExists{upquote.sty}{\usepackage{upquote}}{}
\IfFileExists{microtype.sty}{% use microtype if available
  \usepackage[]{microtype}
  \UseMicrotypeSet[protrusion]{basicmath} % disable protrusion for tt fonts
}{}
\makeatletter
\@ifundefined{KOMAClassName}{% if non-KOMA class
  \IfFileExists{parskip.sty}{%
    \usepackage{parskip}
  }{% else
    \setlength{\parindent}{0pt}
    \setlength{\parskip}{6pt plus 2pt minus 1pt}}
}{% if KOMA class
  \KOMAoptions{parskip=half}}
\makeatother
\usepackage{xcolor}
\usepackage[margin=0.85in]{geometry}
\setlength{\emergencystretch}{3em} % prevent overfull lines
\providecommand{\tightlist}{%
  \setlength{\itemsep}{0pt}\setlength{\parskip}{0pt}}
\setcounter{secnumdepth}{-\maxdimen} % remove section numbering
\usepackage{amsmath,amssymb,amsthm}
\usepackage{longtable}
\ifLuaTeX
  \usepackage{selnolig}  % disable illegal ligatures
\fi
\usepackage{bookmark}
\IfFileExists{xurl.sty}{\usepackage{xurl}}{} % add URL line breaks if available
\urlstyle{same}
\hypersetup{
  pdftitle={Big references, fixed parts, and an original-carrier equation on point blowups},
  pdfauthor={Independent reconstruction for original P16.5},
  hidelinks,
  pdfcreator={LaTeX via pandoc}}

\title{Big references, fixed parts, and an original-carrier equation on
point blowups}
\author{Independent reconstruction for original P16.5}
\date{6 October 2026}

\begin{document}
\maketitle

\section{Abstract and precise status}\label{abstract-and-precise-status}

Original P16.5 remains \textbf{PARTIAL}, candidate=false. This paper
checks the published numerical envelope input in the entire big rational
domain, including its explicit decreasing multiplier-ideal Fubini--Study
net. It proves that the entire genuine cone of a big nonnef class is
disjoint from the published ambient signed finite-energy space; this
obstructs direct use of that product construction, without ruling out a
separate relative theory. It reconstructs the genuine cone translation
for every big rational bundle on a smooth projective surface. That cone
result alone supplies no surface equation. We do not adopt the supplied
unreviewed nef-and-big theorem.

We also prove a narrower equation theorem without that unreviewed
theorem: for every point blowup of projective space of dimension at
least two, and every actual rational bundle \(L=dH+aE\) with \(d,a>0\),
the entire originally defined numerical cone has intrinsic minimal
reference \(a\log\|s_E\|\). Its reference-relative finite-energy space,
genuine metric operator of volume mass \(d^n\), exact original-carrier
target condition, weighted decreasing-net limits, comparison, both
continuous-direction derivatives, actual attainment, and uniqueness
modulo constants are constructed. This is a fixed-part family in
arbitrary dimension, not the entire big domain. The only equation
premise is the task-granted actual-ample G0 theorem, used on an actual
ample bundle on projective space. All functions, outputs, and continuous
observables of the resulting equation remain on the blowup.

At zero volume we establish a constants-only numerical-zero endpoint and
a zero-top-product example with nonuniqueness. We supply no general
nonbig positive-mass replacement. The general big nonnef relative
operator and equation, beyond the stated fixed-part range, remain
unproved. These are unverified proof drafts; source retrieval, rational
diagnostics, and compilation do not certify them.

\section{1. Original data and the sole equation
premise}\label{original-data-and-the-sole-equation-premise}

The ground field \(k\) is algebraically closed, trivially valued, of
characteristic zero, and of arbitrary cardinality. The original carrier
is \(C=X^{\mathrm{an}}\) for a smooth integral projective \(k\)-variety.
A rational line bundle means an actual element of
\(\operatorname{Pic}(X)_{\mathbb Q}\); tensor identities are retained
and cleared by powers when sections are used. Measures are positive
regular Borel Radon measures on the compact Hausdorff carrier. The
generic trivial valuation is \(v_0\).

For \(\theta=c_1(L)\), use the literal original numerical cone:
\(u\in\operatorname{PSH}(X,\theta)\) means that \(u\) is generically
finite, upper semicontinuous, and is the all-point decreasing limit of
one net \[
 u_i=m_i^{-1}\max_j\{\log\|s_{ij}\|_{\mathrm{triv}}+\lambda_{ij}\},
 \quad 0\ne s_{ij}\in H^0(X,m_iB_i),\quad c_1(B_i)\longrightarrow\theta.       \tag{1}
\] Here \(B_i\) are actual rational bundles and weights are real. Common
zeros have value \(-\infty\). The canonical trivial metric evaluates in
a norm-one frame at the valuation-ring center. A generating section
family has zero-weight maximum exactly zero at every analytic point.
Consequently multiplying sections by generators of an actual ample
difference changes their bundle but preserves their scalar function
exactly. Tensoring section families proves addition of cones. Constants
change weights.

Every cone in use embeds into a sufficiently large actual ample cone: if
\(M-L\) is ample, then \(M-B_i\) is ample on a tail, and the preceding
section operation embeds that tail into \(\operatorname{PSH}(X,M)\). The
supplied \textbf{G0}, reproduced exactly in
\texttt{evidence/supplied-14.txt}, is the sole equation grant. Its
stated domain is actual ample rational bundles on geometrically
unibranch integral projective varieties over the fields above. In that
domain we use its fixed-bundle decreasing approximation, normalized
divisorial compactness, all-point divisorial determination and order,
genuine envelopes, raw mixed metric products and pairing, energy,
weighted decreasing-net continuity, comparison/domination, both outer
signs, exact finite-energy targets, target-integral continuity on energy
sublevels, and attained equation. Smooth integral varieties used here
satisfy its geometric hypotheses.

In particular, common ample embedding gives \(u(v_0)=\sup_Cu\), finite
values at every divisorial valuation, and all-point determination of
every genuine function used below. We never infer that an unbounded
nonample function belongs to the published signed finite-energy space
\(\widetilde{\mathcal E}^{1}\).

For bounded nonample nef inputs only, we also use the actual signed
pairing of BJ, arXiv:1801.08229v4, Theorems 7.14 and 7.18 and
Proposition 7.19. Its full pertinent proof is supplied at original
literal LF3880--4057. If \(P\) is nef and \(b\) is bounded \(P\)-psh,
then \(b\) is psh in the actual ample class \(P+T\) for an ample \(T\),
and its two constant bounds make its ample energy finite. Thus its
signed-domain membership is proved before that pairing is used.
Positivity for the nef products follows either from the published
theorem or from the finite polynomial limit of positive ample products
for \(P+\epsilon T\). Their mass is the ordinary nef intersection. This
bounded statement has no unbounded nonample inverse conclusion.

The indispensable bounded construction has no inverse-measure premise.
The actual ample positive pairs span
\(N^1(X)\times\widetilde{\mathcal E}^1\): enlarge the ample class for a
given finite-energy difference, and write any numerical class as a
difference of ample classes. Symmetric multilinearity extends their
actual energy pairing uniquely to that vector space. Its restriction to
PL pairs agrees with the model intersection pairing by the same spanning
argument on the PL subspace. For nef bounded inputs, perturb by
\(\epsilon T\) to deduce pairing monotonicity from actual ample
monotonicity. For each PL test \(f\), assign its pairing with the
product. If \(f\ge0\), write \(f\) as a difference of two bounded
ample-psh PL functions; monotonicity makes this assignment positive. Its
value on one is the ordinary intersection mass, so it is bounded in sup
norm. Original PL density and Riesz construct a unique Radon product on
all original continuous tests. Common-ample decreasing pairing
continuity gives bounded product continuity. This reconstructs the
pertinent proof of the published signed extension within the membership
domain verified above.

\section{2. The complete big domain: the reference input and its exact
limit}\label{the-complete-big-domain-the-reference-input-and-its-exact-limit}

Let \(X\) now be any smooth integral projective variety in the retained
field domain and \(L\) any big actual rational bundle, with no nef,
dimension, countability, support, or finite-generation restriction.
Boucksom--Jonsson's addendum, arXiv:2206.07183v2 (arXiv date 6 September
2023, manuscript date 7 September 2023), Theorem A and Corollaries 2.1
and 2.2, are external published premises. The complete pertinent
statement/proof units were freshly opened. They give the genuine minimal
reference \[V_\theta=P_\theta(0),\qquad V_\theta(v_0)=0,\] and
normalized divisorial compactness. These results are not an equation
grant. For an upper bounded usc obstacle with a genuine minorant, its
usc envelope is a genuine competitor, so the envelope itself is genuine.
Normalized potentials satisfy \(u\le V_\theta\le0\). Thus the reference
is intrinsic to the original cone and canonical metric convention.

We check the crucial big rational approximation directly. Fix a rational
PL obstacle \(f\). On the source's resolved, equivariant, integrally
closed test configuration let \(\mathcal L\) represent \(f\), let
\(\mathfrak a_m\) be the base ideals of \(m\mathcal L\), and let
\(\mathfrak b_m\) be their asymptotic multiplier ideals. Retain the
primary inputs \(\mathfrak a_m\subset\mathfrak b_m\),
\(\mathfrak b_{ml}\subset\mathfrak b_m^l\), and equivariant uniform
generation of \((m\mathcal L+\mathcal A)\otimes\mathfrak b_m\) for one
actual ample \(\mathcal A\). Its restriction is an actual ample bundle
\(A\) on \(X\). The vertical shift of \(\mathcal A\) is chosen so its PL
scalar \(f_A\) is nonnegative.

The source's finite section construction, including Lemma 1.7's section
criterion, gives literal generically finite FS functions
\[g_m=f+m^{-1}\phi_{\mathfrak b_m}+m^{-1}f_A
       \in\mathcal H^{\mathrm{gf}}_{\mathbb Q}(L+A/m).                  \tag{2}\]
Index sufficiently divisible integers by divisibility. Subadditivity
gives
\[(ml)^{-1}\phi_{\mathfrak b_{ml}}\le m^{-1}\phi_{\mathfrak b_m},\] at
all points, including poles. Also \(f_A/(ml)\le f_A/m\). Therefore this
is \textbf{one decreasing original FS net} whose actual numerical
classes tend to \(c_1(L)\). Its generic value has a finite lower bound:
a nonzero big section, shifted below the bounded obstacle, is a
competitor for \(Q_L(f)\). The base-ideal section criterion and
inclusion give \[Q_L(f)\le g_m\le f+f_A/m.\] Its endpoint is generically
finite and usc and belongs to the cone by definition (1), without a
flattening lemma. It is at most \(f\). The equality \(Q_L(f)=P_L(f)\)
off the augmented base locus, proved in Lemma 1.5, puts every psh
competitor below this endpoint at every full-support divisorial
valuation. Common ample all-point order then proves that the endpoint is
precisely \(P_L(f)\). For \(f=0\) this gives a direct original defining
net for \(V_\theta\).

This check retains the source's resolution and multiplier-ideal theorems
as published recursive background. The complete envelope and compactness
results still have the original psh-topology foundations of the addendum
and BJ; we do not claim a new recursive proof of that entire theory. In
particular our direct check of (2) does not certify every possible joint
closure statement. It establishes the specific single-net membership
needed for this reference. This prevents a defective abstract flattening
argument from being used as its substitute.

The envelopes and reference floors \(\max(u,V_\theta-k)\) exist in the
full big domain. There is, at this stage, \textbf{no constructed
relative energy or volume-mass operator} on those floors. Sections 7--9
below construct a relative theory only when a separately proved
fixed-part cone translation supplies actual bounded residual products.

There is a stronger, general domain obstruction. Under these reference
premises, if \(\theta\) is nonnef then
\[\operatorname{PSH}(\theta)\cap\widetilde{\mathcal E}^{1}=\varnothing. \tag{2a}\]
To prove this, the scaling action \(t\cdot u(v)=t u(t^{-1}v)\) preserves
(1), merely replacing each FS weight \(\lambda\) by \(t\lambda\). It
preserves competitors under the zero obstacle and is invertible, so
\(t\cdot V_\theta=V_\theta\). Thus the reference is homogeneous. Embed
it in an actual ample cone \(\Omega\) with \(\Omega-L\) ample. If it had
finite \(\Omega\)-energy, the model base-change identity, passed through
actual ample decreasing energy continuity, would give
\[E_{\Omega,\mathrm{raw}}(V_\theta)=E_{\Omega,\mathrm{raw}}(2\cdot V_\theta)
      =2E_{\Omega,\mathrm{raw}}(V_\theta).\] Its normalized energy would
be zero. The actual ample \(I,J\) inequality and all-point \(I\)
separation in G0 then imply \(V_\theta=0\). That would make zero a
bounded \(\theta\)-psh potential and force \(\theta\) nef by the literal
definition (bounded endpoints force finite generating FS stages and
hence nef class limits). This contradiction proves
\(E_\Omega(V_\theta)=-\infty\).

Every normalized \(u\in\operatorname{PSH}(\theta)\) satisfies
\(u\le V_\theta\). If \(u\) had finite \(\Omega\)-energy, upward
finite-energy closure would give finite energy to \(V_\theta\), another
contradiction. Finally BJ Lemma 7.13, applied at the actual ample
\(\Omega\), identifies
\(\widetilde{\mathcal E}^{1}\cap\operatorname{PSH}(\Omega)\) with
\(\mathcal E^1(\Omega)\). This proves (2a). All relative floors are
genuine functions in the same cone and hence are excluded too. This is
an obstruction to applying the published signed pairing directly, not a
nonexistence theorem for a separately constructed relative operator. It
applies to the entire retained big nonnef domain, independently of a
finite fixed divisor decomposition.

\section{3. A genuine nef intersection
lemma}\label{a-genuine-nef-intersection-lemma}

This lemma is used for the surface cone result, not as a nonample
equation premise. For a nef actual rational bundle \(P\) and ample
\(T\), one has
\[\operatorname{PSH}(P)=\bigcap_{\epsilon\in\mathbb Q_{>0}}
             \operatorname{PSH}(P+\epsilon T).                       \tag{3}\]
The forward inclusion follows by multiplying a defining section tail by
generators of the ample difference. Conversely take a generically finite
function \(v\) in every approaching ample cone. In each fixed ample
bundle use G0 to choose finite continuous FS majorants decreasing to
\(v\). Add strictly positive shifts \(\delta\), and order triples by
decreasing \(\epsilon\), decreasing \(\delta\), and decreasing shifted
scalar.

Given two triples, choose a smaller positive rational \(\epsilon_3\) and
\(0<\delta_3<\min(\delta_1,\delta_2)\). The finite continuous old
minimum is strictly above \(v+\delta_3\). The open sets where the
fixed-\(\epsilon_3\) continuous majorants plus \(\delta_3\) are below
that minimum form a directed cover of compact \(C\). One late majorant
is below the old minimum everywhere. This proves directedness, even
where \(v=-\infty\). Requiring arbitrarily small parameters and
refinement of each fixed-class majorant shows that the resulting single
original FS net has endpoint \(v\) and class limit \(P\). This proves
(3).

For \(\epsilon\le1\), the normalized approaching cones are compact
closed subsets of one fixed larger ample cone, in the same divisorial
Hausdorff topology. Their intersection is therefore compact. This
establishes normalized nef compactness and actual function membership
needed below, without importing a nef equation or a general flattening
lemma. Zero is genuine by zero-weight generating sections of
\(P+\epsilon T\).

\section{4. Every big surface bundle: actual fixed sections and cone
translation}\label{every-big-surface-bundle-actual-fixed-sections-and-cone-translation}

Suppose \(\dim X=2\) and \(L\) is an arbitrary big actual rational
bundle. Choose a section of a divisible power and divide its effective
divisor by that power. The surface Zariski theorem gives
\[L=P+N,\quad N=\sum_{j=1}^r a_jC_j\ge0,\] where \(P=L-\mathcal O(N)\)
is an actual nef rational bundle, \(M=(C_i\cdot C_j)\) is negative
definite if \(N\ne0\), and \(P\cdot C_i=0\). We import this algebraic
theorem, with its complete pertinent proof freshly checked in Bauer,
arXiv:0712.1576v1, and check all applications below. It does not grant
any analytic equation.

First \(M^{-1}\) is entrywise nonpositive. Indeed \(A=-M\) is positive
definite with nonpositive off-diagonal entries. If \(y\ge0\) and
\(x=A^{-1}y=x_+-x_-\), disjoint nonnegative parts give
\[x_-^Ty=x_-^TAx_+-x_-^TAx_-<0\] if \(x_-\ne0\), a contradiction. Thus
\(A^{-1}y\ge0\).

For any effective rational divisor \(F\) of numerical class \(B\), write
\(F=\sum d_jC_j+R\) with none of the \(C_j\) a component of \(R\). Put
\(e_i=(B-L)\cdot C_i\). Distinct effective curves have nonnegative
intersections, so
\[M(d-a)=e-(R\cdot C_i)_i\le e,\qquad d-a\ge M^{-1}e.                  \tag{4}\]
If \(\eta\ge\max_i|e_i|\) and \(c_j=\sum_i|(M^{-1})_{ji}|\), every such
\(F\) contains \[N_\eta=\sum_j\max(a_j-c_j\eta,0)C_j.\] This is uniform
over every section in every degree of \(B\), and \(N_\eta\to N\) as
\(B\to L\). At \(B=L\), it gives actual section identities for all
sufficiently divisible \(m\):
\[H^0(X,mL)=s_{mN}H^0(X,mP).                                        \tag{5}\]
Division in (5) is division of sections by a fixed Cartier section, not
subtraction of infinite scalar values. Consequently \(P\) is
section-big. The same matrix argument applied in both directions to two
effective representatives shows that \(N\) is intrinsic: every other
representative contains \(N\), and its residual effective nef subdivisor
forces its Zariski negative part to be at most \(N\) by maximality;
reversing the roles gives equality. Numerically trivial twists remain in
the actual bundle \(P\).

For clarity, the needed nef volume equality can be proved independently
of the analytic drafts. For integral nef \(P\) on a smooth projective
\(n\)-fold, choose smooth very ample \(B\) with \(B-K_X\) ample (for
\(n=1\), a reduced finite divisor). Kodaira vanishing and the
intersection Euler polynomial give
\[h^0(X,mP+B)=P^n m^n/n!+O(m^{n-1}).\] The restriction has
\(O(m^{n-1})\) sections: on the integral divisor choose very ample \(T\)
with \(T-P|_B\) ample; a section of \(m(T-P|_B)\) embeds its restriction
sections into those of \(mT+B|_B\), bounded by the Hilbert polynomial.
The actual Cartier restriction sequence sandwiches \(h^0(X,mP)\) between
the displayed polynomial and that polynomial minus \(O(m^{n-1})\). Thus
\(\operatorname{vol}(P)=P^n\). For rational bundles clear denominators.
Kodaira vanishing applies over arbitrary characteristic-zero \(k\) by
descending each finite presentation to a finitely generated subfield,
embedding its algebraic closure into \(\mathbb C\), and using flat
cohomology base change. This imposes no cardinality condition on \(k\).
These named classical algebraic theorems remain recursive background. By
(5), \(\operatorname{vol}(L)=P^2>0\).

Write \(\sigma_N=\log\|s_N\|_{\mathrm{triv}}\). We now prove the entire
cone equality
\[\operatorname{PSH}(L)=\sigma_N+\operatorname{PSH}(P).               \tag{6}\]
Take the original decreasing FS net \(u_i\) for
\(u\in\operatorname{PSH}(L)\). Formula (4) factors every section in its
stage by one rational \(N_i\le N\), with \(N_i\to N\). Taking a common
further power clears all denominators and leaves the scalar unchanged.
The residual scalar \(v_i\) is an actual FS function in \(B_i-N_i\),
with generic value \(b_i=u_i(v_0)\downarrow b=u(v_0)\). The residual net
need not decrease.

For each fixed rational \(\epsilon>0\), \(P+\epsilon T-(B_i-N_i)\) is
ample on a tail. Normalize \(v_i-b_i\) and embed into one larger ample
cone. Compactness supplies a cluster. At each original divisorial \(w\),
all divisor scalars are finite, and its coordinate is uniquely
prescribed by \[v_i(w)-b_i\longrightarrow u(w)-\sigma_N(w)-b.\] For
every \(\epsilon\), the original tail is in the closed compact
normalized \(P+\epsilon T\) cone. The cluster belongs to their
intersection, hence to \(\operatorname{PSH}(P)\) by (3). Restore \(b\).
Multiplying its genuine defining net by \(s_N\) produces a decreasing
original net in classes tending to \(L\). Its endpoint agrees with \(u\)
at every divisorial valuation; common ample all-point determination
proves equality everywhere, including \(N^{\mathrm{an}}\). Conversely
the same multiplication proves the other inclusion in (6). Uniqueness of
the residual follows at divisorial coordinates and then everywhere. Both
maps are continuous in those coordinates and reflect order.

The normalized residual maximum principle now gives exactly
\(V_L=\sigma_N\). This checks genuine membership, actual tensors,
coordinate limits, and all-point identity in the supplied surface
proposal. It does \textbf{not} independently establish its relative
operator or equation for arbitrary nef and big \(P\). Defining an energy
formally as \(E_P\) would still require the entire actual intrinsic nef
theorem. That theorem is an unreviewed draft here. The following
independent equation uses a different, narrower residual whose equation
lies within G0.

\section{5. A nonnef fixed-part family in arbitrary
dimension}\label{a-nonnef-fixed-part-family-in-arbitrary-dimension}

Let \(Y=\mathbb P^n_k\), \(n\ge2\), let \(p\in Y(k)\), and let
\[q:X=\operatorname{Bl}_pY\longrightarrow Y.\] Write
\(H=q^*\mathcal O_Y(1)\) and \(E\) for the exceptional Cartier divisor.
Retain \textbf{all} \(d,a\in\mathbb Q_{>0}\) and actual \(L=dH+aE\). Put
\(A=\mathcal O_Y(d)\), \(P=q^*A\), and
\(\sigma=a\log\|s_E\|_{\mathrm{triv}}\).

We record the elementary algebra. Every rational bundle on \(X\) is
uniquely \(bH+cE\). A divisor of a rational section pushes to a Cartier
divisor on the smooth factorial \(\mathbb P^n\); subtract its pullback
and only \(E\) remains. Linear independence follows by intersection with
a line in an exceptional fiber (\(E\) has degree \(-1\), \(H\) degree
zero), and with a curve avoiding \(p\). This proves the actual Picard
statement, not just numerical generation, and gives the same two
coordinates in \(N^1(X)\). There is no ignored numerically trivial twist
on this \(X\).

If \(b,c\ge0\) are integral, restriction off \(E\) and Hartogs on normal
\(Y\) prove
\[H^0(X,bH+cE)=s_E^c q^*H^0(Y,\mathcal O_Y(b)).                       \tag{7}\]
Indeed the restriction of the section divided by \(s_E^c\) is a section
of \(\mathcal O_Y(b)\) off a codimension-at-least-two point. In a local
polynomial UFD it cannot have a denominator: each irreducible
denominator would have a codimension-one generic pole distinct from
\(p\). Thus it extends, and equality off \(E\) identifies the global
sections. This proves both inclusions. Rational statements follow after
clearing powers.

Consequently
\[h^0(X,mL)=\binom{md+n}{n},\qquad \operatorname{vol}(L)=d^n=:V>0,     \tag{8}\]
for divisible \(m\). The degree of \(L\) on an exceptional line is
\(-a\), so every bundle in this family is nonnef. Its ordinary top
intersection is \[L^n=d^n+(-1)^{n-1}a^n,\] since \(H|_E\) is trivial and
\(E|_E=\mathcal O_{\mathbb P^{n-1}}(-1)\). It is generally not the
equation's mass.

\textbf{Cone theorem.} The entire originally defined cone is
\[\operatorname{PSH}(X,L)=\sigma+q^*\operatorname{PSH}(Y,A).          \tag{9}\]
To prove it, write a defining stage in \(B_i=d_iH+a_iE\). Eventually
\(d_i,a_i>0\). Formula (7) factors the stage exactly as
\[u_i=a_i\log\|s_E\|+q^*f_i,\qquad f_i\in\mathcal H^{\mathrm{gf}}(\mathcal O_Y(d_i)).\]
Here \(d_i\to d\), \(a_i\to a\), and
\(f_i(v_{0,Y})=u_i(v_{0,X})=:b_i\downarrow b\). Choose
\(\Omega=\mathcal O_Y(d+1)\). The normalized residuals \(f_i-b_i\) embed
in its normalized compact cone. At every base divisorial valuation,
identified with the common-field valuation upstairs, their coordinates
converge to \(u(w)-\sigma(w)-b\). A cluster \(v\) has exactly those
coordinates. For each rational \(\epsilon>0\) its tail lies in the
closed compact \((d+\epsilon)\mathcal O_Y(1)\) cone, so does \(v\).

Because \(v\le0\), the functions \(d(d+\epsilon)^{-1}v\) are genuine
\(A\)-psh and decrease to \(v\) as \(\epsilon\downarrow0\). G0's fixed
ample decreasing closure proves \(v\in\operatorname{PSH}(Y,A)\)
directly. Restore \(b\). Pull a fixed-bundle G0 net for \(v+b\) back and
multiply by \(s_{aE}\). It is a single decreasing original FS net in
\(L\). Its endpoint agrees with \(u\) at every divisorial valuation, so
agrees everywhere in a common ample cone on \(X\). This proves (9). The
converse is the same actual section multiplication. Residuals need not
be decreasing before their limit is identified.

The map in (9) is an order isomorphism and a divisorial homeomorphism.
To reflect order, subtract the finite divisor scalar at full-support
divisorial valuations and then apply G0 all-point order on \(Y\). To
prove continuity, those same coordinate equalities suffice. We have
\[V_L=\sigma.\] The reference and residual \(R(u)=q^*v\) are intrinsic:
\(R(u)\) is the unique genuine \(P\)-potential extending \(u-V_L\) off
\(E^{\mathrm{an}}\). It is defined on all of \(X^{\mathrm{an}}\),
including the common pole set of \(u,V_L\), by this proved unique
extension, rather than by subtracting two infinite values.

\section{6. A lifting lemma for every original continuous
observable}\label{a-lifting-lemma-for-every-original-continuous-observable}

We need a topological fact to keep the original carrier. Suppose
\(q:C\to D\) is a continuous surjection of compact Hausdorff spaces and
is a homeomorphism between \(C\setminus F\) and \(D\setminus Z\), where
\(F=q^{-1}Z\) and \(Z\) is closed. A positive finite Radon measure
\(\nu\) on \(D\) with \(\nu(Z)=0\) has exactly one Radon measure
\(\Lambda\nu\) on \(C\) with \(q_*\Lambda\nu=\nu\) and
\(\Lambda\nu(F)=0\). Restrict to the open complement, use its actual
inverse homeomorphism, and extend by zero. Radon regularity follows by
choosing a compact \(K\subset D\setminus Z\) with arbitrarily small
complementary mass and pulling it back. This proves existence and
uniqueness as measures, not an inverse metric operator.

If \(\nu_i\to\nu\) weakly, all \(\nu_i(Z)=0\), and \(\nu(Z)=0\), then
\(\Lambda\nu_i\to\Lambda\nu\) on \textbf{every} \(f\in C(C)\). Let
\(\widetilde f=f\circ q^{-1}\) on \(D\setminus Z\). For any
\(\epsilon>0\), regularity and Urysohn give \(0\le\chi\le1\), equal to
one on a neighborhood of \(Z\), with \(\int\chi\,d\nu<\epsilon\). The
function \((1-\chi)\widetilde f\), extended by zero, is continuous on
\(D\). Its integrals converge and its replacement error is at most
\(\|f\|\int\chi\,d\nu_i\), eventually small. The same holds at the
limit. This proves the full net assertion. For signed weights with
positive weighted measures converging weakly and no limit mass on \(Z\),
apply the same proof to those positive measures; constants and
differences then restore signed weights.

For our blowup, \(Z=\{p\}\), \(F=E^{\mathrm{an}}\), and the inverse is
the actual isomorphism off these algebraic loci. No retraction,
skeleton, or arbitrary choice of point in a fiber occurs. This lemma
will identify limits of preconstructed metric products; it does not
define them by prescribed targets.

\section{7. Polar exclusion, actual bounded products, and relative
energy}\label{polar-exclusion-actual-bounded-products-and-relative-energy}

We prove the one polar input we need inside the ample domain. The
degree-\(d\) FS scalar made from the hyperplanes through \(p\) is a
normalized genuine \(A\)-psh function \(h\le0\), with poles exactly
\(p\). Choose \(0<\gamma<1/(n+1)\), and set \(b=1-(1-h)^\gamma\). This
is genuine \(A\)-psh: its convex increasing scalar map, with slopes
between zero and one, is the decreasing limit of finite piecewise-affine
chord majorants with constant left tails. Each majorant is a finite
maximum of slopes times \(h\) plus constants, retaining unused ample
weight on zero.

Here is the finite-energy check. With the ample probability notation
only after \(V>0\), put
\(\operatorname{Cap}(B)=\sup_{-1\le c\le0}\operatorname{MA}_A(c)(B)\).
Bounded mixed integration by parts, telescoped from the canonical
measure \(\delta_{v_0}\), gives
\(\int(-h_s)\operatorname{MA}_A(c)\le n\) for \(h_s=\max(h,-s)\). Each
difference is the integral of \(-c\in[0,1]\) against a difference of
positive mixed products of unit mass. Countable floor convergence gives
\(\operatorname{Cap}\{h\le-t\}\le n/t\).

For any normalized genuine \(z\), locality identifies the measures of
its \(s\)- and \(t\)-floors on \(\{z>-t\}\), for \(s\ge t\). Equal mass
identifies their complementary masses. The positive expansion of
\(t^n\operatorname{MA}(z_t/t)\) dominates \(\operatorname{MA}(z_t)\) for
\(t\ge1\). Layer cake and the energy supporting inequality yield
\[-E_A(z_s)/V\le1+\int_1^\infty t^n\operatorname{Cap}\{z\le-t\}\,dt.\]
For \(z=b\), its deep set is \(\{h\le1-(1+t)^{1/\gamma}\}\), so the
integrand is \(O(t^{n-1/\gamma})\), integrable by our explicit
conservative choice of \(\gamma\). Thus \(b\in\mathcal E^1(Y,A)\) and
\(b(p)=-\infty\). This proof does not rely on a stronger printed
attenuation exponent. Every ample finite-energy target, and every top or
mixed finite-energy output, integrates \(b\) by G0. They therefore give
\(p\) zero mass.

For bounded \(v_j\in\operatorname{PSH}(Y,A)\), construct on \(X\) the
bounded nef metric product \[\mathcal M_P(q^*v_1,\ldots,q^*v_n)
       =\bigwedge_{j=1}^n(P+dd^c q^*v_j)\] in the verified signed domain
of Section 1. At a common compactified model this is the actual residual
intersection measure
\[\sum_F b_F\prod_{j=1}^n(\mathcal P_j+D_j)|_F\,\delta_{w_F},
  \quad b_F=\operatorname{ord}_F(t),\quad
  w_F=b_F^{-1}\operatorname{ord}_F|_{k(X)}.                            \tag{10}\]
The residual pairs are pulled from the base metric models. Projection of
the full fundamental central cycle proves that its pushforward equals
the actual ample base mixed product. Its mass is \(A^n=V\). A central
component mapping to smaller dimension has zero pulled \(n\)-fold
degree; no central multiplicity is omitted. The pairing identity for
every pulled PL test is equally the projection formula. For bounded
noncontinuous \(v_j\), take fixed-base ample decreasing PL nets; pull
them back and use bounded decreasing signed-pairing continuity on \(X\),
already verified by bounded membership in an actual common ample cone.
PL density on \(Y\) gives the pushforward identity for every base
continuous test.

The base product gives \(p\) zero mass. Thus the positive original
product gives \(E^{\mathrm{an}}\) zero mass and is identified by Section
6 with its \textbf{unique} extension from the actual isomorphism locus.
For continuous inputs this is also the local metric product in the
bundle identity \(L=P+aE\): in the nonvanishing \(s_{aE}\) frame off
\(E\), the \(L\)-metric has the residual \(P\)-weight. The prescribed
singular horizontal divisor has been removed before taking the product.
We retain this zero-mass extension at the reference poles; this is a
reference-relative nonpluripolar product.

For \(u=\sigma+q^*v\) with bounded \(v\), define raw relative energy by
the actual residual pairing
\[E_{L,V_L}(u)=\frac1{n+1}\sum_{j=0}^n\int R(u)\,
  \mathcal M_P(\underbrace{R(u),\ldots,R(u)}_j,
                    \underbrace{0,\ldots,0}_{n-j})=E_{A,\mathrm{raw}}(v). \tag{11}\]
The equality follows from the model projection just proved and bounded
decreasing limits. The zero reference has energy zero. Define
\(u_k=\max(u,V_L-k)=\sigma+q^*\max(v,-k)\); equality includes the pole
set, through the unique residual. Put
\[E_{L,V_L}(u)=\inf_k E_{L,V_L}(u_k),\qquad
\mathcal E^1(L,V_L)=\{u\in\operatorname{PSH}(L):E_{L,V_L}(u)>-\infty\}. \tag{12}\]
It is exactly \(\sigma+q^*\mathcal E^1(Y,A)\) by (11) and ample floor
continuity. It is the entire actual intrinsic reference-floor space in
this family, not a declared ambient signed space. Constants shift raw
energy by \(Vc\). The reference is target-independent and prescribed in
every floor and energy.

\section{8. The original unbounded operator and weighted
limits}\label{the-original-unbounded-operator-and-weighted-limits}

For \(v\in\mathcal E^1(Y,A)\), take the actual bounded original products
of its floors from Section 7. Their pushforwards converge to the already
constructed ample metric product \(M_{A,\mathrm{raw}}(v)\). The endpoint
and every stage give \(p\) zero mass by Section 7. Section 6 proves that
those \textbf{original metric floor measures} converge on every
continuous test of \(X\). Define their limit as
\[M_{L,V_L}(u)=\lim_{k\to\infty}\mathcal M_P(q^*v_k).
                                                                    \tag{13}\]
It has mass \(V=\operatorname{vol}(L)\) and no mass on
\(E^{\mathrm{an}}\). It is identified, after this product construction,
with \(\Lambda M_{A,\mathrm{raw}}(v)\). The latter is an identity, not
an operator defined by choosing a measure to solve a target. Changing
cofinal real floor levels or a bounded reference does not change the
limit, since the base energy theory and Section 6 give the same base
operator and same unique extension. Auxiliary ample choices were removed
at each bounded stage by uniqueness of the signed metric pairing.

Construct mixed relative products the same way from bounded residual
mixed floors. G0's base mixed continuity and Section 6 identify their
positive full-mass limits, independent of all cofinal floor choices. If
\(u_{j,i}\downarrow u_j\) are original genuine nets with finite-energy
endpoints, their residual base functions decrease all-point: order is
reflected by (9), and their pointwise infimum equals the required
endpoint on all divisorial coordinates and hence everywhere by fixed
ample determination. Generic values remain finite. The base mixed net
converges by G0, and Section 6 proves convergence of its actual original
products against all \(C(X^{\mathrm{an}})\).

Every intrinsic residual test \(R(z)=q^*w\), \(w\in\mathcal E^1(Y,A)\),
is integrable for every such top and mixed product, by the base identity
and G0. For simultaneous decreasing residual weights
\(w_i\downarrow w\), translate below zero on a tail. The positive base
weighted measures \((-w_i)M_A(v_{1,i},\ldots,v_{n,i})\) converge weakly
by G0 and have a finite endpoint, which gives \(p\) zero mass. Section 6
applies directly to these positive weighted measures. This proves the
original integrable weighted net theorem, with the whole original
continuous test space. Restore constants to obtain signed weights. No
assertion of arbitrary-Borel-net monotone convergence is used.

The theorem also handles an actual original psh weight
\(z_i\downarrow z\) whenever \(z\) is integrable for the limiting
product. Write \(z_i=\sigma+q^*w_i\) using (9). The ideal of \(p\)
satisfies \(\mathcal I_p\mathcal O_X=\mathcal O_X(-E)\), so
\(q^*\log|\mathcal I_p|=\log\|s_E\|\) at every point. Thus
\(z_i=q^*(a\log|\mathcal I_p|+w_i)\), where the functions in parentheses
are decreasing psh weights in the actual ample class
\((d+a)\mathcal O_Y(1)\). They need not be finite energy. Apply the
actual-ample integrable-weight theorem (BJ Proposition 7.21, with the
corrected pairing sign, or G0's corresponding weighted statement) to
these weights and the residual \(A\)-products. Endpoint integrability is
exactly the original hypothesis. Their translated negative weighted
measures converge positively, and give \(p\) zero mass, so Section 6
again proves the full original weighted limit. Absolute integrability of
\(\sigma\) is not added to the relative target definition (16); this
additional weighted assertion is explicitly conditional on the stated
endpoint integrability.

All residual operators agree with local original \(L\)-metric products
off the reference poles, and their limit has no mass there. This proves
the claimed metric interpretation and descent within the stated family.
It is not a construction for a general movable b-divisor or a general
big class.

\section{9. Original variational inputs and all continuous
directions}\label{original-variational-inputs-and-all-continuous-directions}

Normalized compactness and compact energy sublevels are transferred
through the \textbf{proved} cone homeomorphism (9), from G0 on the
actual ample \(A\). The relative energy is concave, monotone,
divisorially usc, and decreasing-net continuous by the exact identities
(11)--(12). These conclusions concern actual original functions. Genuine
original envelopes for usc obstacles with a minorant follow from the
published big envelope result in Section 2. In particular
\[q_t=P_L(u+tf),\qquad u-|t|\|f\|\le q_t\le u+|t|\|f\|,\] is actual
relative finite energy for every \(f\in C(X^{\mathrm{an}})\).

We must not assume that \(f\) descends continuously through \(q\) at
\(p\). Put \(\widetilde f=f\circ q^{-1}\) on
\(Y^{\mathrm{an}}\setminus\{p\}\), and \(F=\|f\|\). Choose continuous
\(0\le\eta\le1\), zero on a neighborhood of \(p\). Then
\[g^-=\eta\widetilde f-F(1-\eta),\qquad
  g^+=\eta\widetilde f+F(1-\eta)\] extend continuously to \(Y\), satisfy
\(g^-\le\widetilde f\le g^+\) off \(p\), and can be chosen with
integrals against \(M_A(v)\) arbitrarily close to
\(\int fM_{L,V_L}(u)\), by regularity and absence of mass at \(p\).

Write \(R(q_t)=q^*v_t\). Base energy concavity and the original obstacle
constraint off \(E\) give, for every real \(t\),
\[E_{L,V_L}(q_t)-E_{L,V_L}(u)
 \le\int(v_t-v)M_A(v)\le t\int fM_{L,V_L}(u).                        \tag{14}\]
The integrals are legitimate: the two residuals are finite energy, and
the discarded algebraic fiber has zero mass. For \(t>0\),
\(\sigma+q^*P_A(v+t g^-)\) is a feasible original competitor. Its base
derivative gives the lower derivative bound \(\int g^-M_A(v)\). For
\(t<0\) use \(g^+\); division reverses the competitor inequality and
(14). Taking cutoffs with vanishing lost mass proves both signs:
\[\left.\frac d{dt}\right|_{0}E_{L,V_L}(P_L(u+tf))
             =\int fM_{L,V_L}(u),\qquad f\in C(X^{\mathrm{an}}).       \tag{15}\]
This supplies all original continuous directions; no quotient of the
observable space is substituted. Here outer relative energy means the
energy of the genuine envelope for these usc obstacles. We do not define
it as the infimum over globally finite continuous majorants of \(u\):
that convention could forget the retained singular reference.
Equivalently it is the supremum of relative energies of actual
competitors under this obstacle; genuine envelope membership and
monotonicity identify the two values. The singularity type is retained
throughout.

Define the intrinsic target dual for positive original Radon \(\mu\) of
raw mass \(V\) by \[D_{L,V_L}(\mu)=\sup_{u\in\mathcal E^1(L,V_L)}
  \left(E_{L,V_L}(u)-\int R(u)\,d\mu\right).                         \tag{16}\]
All tests \(R(u)\) are the uniquely extended original residual functions
from Section 5. Exact cone equality gives
\[D_{L,V_L}(\mu)=D_{A,\mathrm{raw}}(q_*\mu).\] G0 makes its finiteness
equivalent to integrability of every original residual finite-energy
test. The test \(q^*b\) from Section 7 forces
\(\mu(E^{\mathrm{an}})=0\), so this is a consequence of the intrinsic
target condition, not an imposed support restriction. Constants make
(16) gauge invariant exactly at mass \(V\).

On compact energy sublevels, target integration is continuous by G0
applied to the actual base measure \(q_*\mu\) and the exact equality of
integrals. Coercivity is likewise the base coercivity in the same
energy. The base attained dual maximizer yields an actual original
attained maximizer by (9) and (12), not a completion object. If \(v\) is
its ample solution, (13) and the unique extension identity give
\[M_{L,V_L}(\sigma+q^*v)=\Lambda(q_*\mu)=\mu.\] Conversely every output
satisfies (16) by base energy concavity. Thus raw mass \(V\) and finite
intrinsic dual are the exact target image.

Equal original output measures have equal base output measures. G0
uniqueness on \(Y\) and injectivity in (9) give equality modulo one real
constant at \textbf{every} original point, including the reference
poles. Comparison follows by pushing the strict inequality set off
\(E\), where it is exactly the base strict inequality set; \(E\) is
null. Domination for a genuine potential against a relative
finite-energy one transfers in the same way, and reflects all-point
order through (9). Mixed and weighted assertions were constructed before
targets in Section 8.

\textbf{Theorem (point-blowup fixed-part equation).} For every retained
field, every \(n\ge2\), every point blowup
\(X=\operatorname{Bl}_p\mathbb P^n\), and every actual rational
\(L=dH+aE\) with \(d,a>0\), (9), (11)--(13), and (16) give a genuine
original-carrier relative finite-energy equation with reference
\(V_L=a\log\|s_E\|\), raw mass \(\operatorname{vol}(L)=d^n\), both
derivatives (15), comparison/domination, common-class mixed and
integrable weighted arbitrary decreasing nets, actual dual attainment,
and uniqueness modulo constants for precisely all targets of that mass
and finite intrinsic dual. Its proofs use G0 only at the actual ample
\(A=\mathcal O_{\mathbb P^n}(d)\). No general nef equation is used. This
narrower theorem does not close original P16.5.

\section{10. Zero-volume results actually
established}\label{zero-volume-results-actually-established}

For every smooth integral projective original \(X\) in the retained
domain and every numerically zero actual rational bundle \(L\), the
genuine numerical cone consists exactly of finite constants. For a
defining net put \(b_i=u_i(v_0)\downarrow b\). For each positive
rational \(\epsilon\), the actual difference \(\epsilon T-B_i\) is ample
on a tail. Thus \(\epsilon^{-1}(u_i-b_i)\) embeds into the normalized
actual ample \(T\) cone. For each divisorial \(w\) its finite lower
bound \(-T_T(w)\) gives \[-\epsilon T_T(w)\le u_i(w)-b_i\le0.\] First
take that tail's limit, then \(\epsilon\downarrow0\). All divisorial
coordinates of \(u\) equal \(b\), hence so do all original values by
common ample determination. Constants are genuine from zero-weight
generating families in classes \(\epsilon T\) tending to zero. Bounded
nef signed products have mass zero and vanish. The relative raw top
energy is identically zero and its floor space is this constants-only
cone. The sole positive target of mass zero is the zero measure. The
equation is attained uniquely modulo constants. There is no probability
normalization. A numerically trivial bundle has section volume zero: a
nonzero section has an effective divisor of zero ample intersection,
hence no divisor, so a cleared power trivializes and has at most one
section.

There is no general zero-mass uniqueness. On
\(X=\mathbb P^1\times\mathbb P^1\) let \(L=p_1^*\mathcal O(1)\). The
task-granted all-point fiber theorem with actual ample base gives fiber
constancy for the entire genuine numerical cone. Pull its defining net
back along the second-coordinate-zero section. The endpoint is
generically finite by that fiber constancy, hence this is a genuine base
defining net. Conversely base fixed-bundle nets pull back. Therefore the
entire cone is \(p_1^*\operatorname{PSH}(\mathbb P^1,\mathcal O(1))\).

Its nef bounded top products are positive of mass \(L^2=0\) and thus
zero; the bounded energy is zero, and all reference floors give zero top
energy and operator on the entire cone. The volume is zero since
\(h^0(mL)=m+1\). The cone has nonconstant bounded functions, for
instance a pullback of \(\max(\log\|z_1\|,-1)\). The zero target has
multiple solutions differing by nonconstant functions. This disproves
retention of top-degree uniqueness for this nonbig class. It does not
refute every lower-degree theory with an extra retained datum. We
establish no such general lower-degree theorem here.

\section{11. Remaining exact original
obligations}\label{remaining-exact-original-obligations}

The full big rational domain now has a checked published minimal
reference and compact normalized genuine cone. For arbitrary big nonnef
\(L\) in higher dimension, we have not constructed positive
reference-relative products on its generally unbounded floors with mass
\(\operatorname{vol}(L)\). We have not established their intrinsic
relative energy, mixed and weighted limits, comparison/domination, both
outer signs, exact target image, original-function attainment, or
all-point uniqueness. A moving model or asymptotic section filtration
would need a target-independent actual metric construction, independence
of its approximations, and a proved original-carrier descent. No finite
Zariski decomposition or fixed birational contraction covers that domain
here.

For arbitrary big surfaces we proved (6), its intrinsic reference, and
volume. The equation still needs a fully independently reconstructed
intrinsic nef-and-big theorem for every possible \(P\); the supplied
proofs are not grants. Our theorem proves only the explicit
residual-ample point-blowup family. We do not state the supplied surface
equation as an unconditional conclusion.

For general nonbig nonzero pseudoeffective classes, the zero top mass
cannot be divided by a probability normalization. A replacement of
positive degree or extra positivity datum and its exact original target
image and uniqueness relation remain unproved. The numerical-zero
theorem and product counterexample are the separately established nonbig
range.

This paper therefore provides new proofs and checked source
applications, and precisely delimited failures, while leaving
\textbf{PARTIAL}, candidate=false for the exact unchanged original
P16.5. The source appendix and ledger distinguish grants, imported
published premises, own deductions, and recursive limitations. No
independent-review verdict is used.

\end{document}
